% arXiv-compatible LaTeX - sample paper
% Compatible with BasicTeX

\documentclass[11pt,a4paper]{article}

% Core packages - all in BasicTeX
\usepackage[utf8]{inputenc}
\usepackage[T1]{fontenc}
\usepackage{lmodern}
\usepackage{microtype}
\usepackage{amsmath,amssymb,amsfonts}
\usepackage[numbers,sort&compress]{natbib}
\bibliographystyle{unsrtnat}
\usepackage{hyperref}
\usepackage{url}
\usepackage{booktabs}
\usepackage{graphicx}
\usepackage[margin=1in]{geometry}
\usepackage{xcolor}

\hypersetup{
  pdfauthor={Cameron Warren},
  pdftitle={OSINT Methods for ICS Security Discovery},
  colorlinks=true,
  linkcolor=blue,
  citecolor=blue
}

\title{Open Source Intelligence Methods for Discovering\\
       Internet-Exposed Industrial Control Systems}
\author{
  Cameron Warren\\
  \texttt{cwarre33@uncc.edu}\\[0.5em]
  Department of Computer Science\\
  University of North Carolina at Charlotte
}
\date{\today}

\begin{document}

\maketitle

\begin{abstract}
Industrial Control Systems (ICS) increasingly connect to public networks, 
creating attack surfaces that threat actors exploit. This paper presents an 
Open Source Intelligence (OSINT) methodology for discovering exposed ICS assets. 
We describe a pipeline comprising automated discovery, service enumeration, and 
vulnerability correlation. Our approach uses targeted scanning for protocols 
including Modbus, BACnet, DNP3, and S7. We discuss ethical frameworks governing 
this research and propose responsible disclosure practices.
\end{abstract}

\section{Introduction}

The convergence of Operational Technology (OT) and Information Technology (IT) 
has expanded the attack surface available to adversaries. Unlike traditional IT, 
ICS were not designed with security as a primary consideration.

\textbf{Contribution.} This work presents: (1) a reproducible OSINT framework 
for ICS asset discovery; (2) analysis of exposure patterns across protocols; 
(3) responsible disclosure recommendations.

\section{Background}

\subsection{ICS Protocols}

Industrial control systems rely on several specialized protocols:

\begin{itemize}
  \item \textbf{Modbus (TCP/502):} No built-in authentication; function codes 
    control read/write operations.
  \item \textbf{BACnet (UDP/47808):} Building automation via Who-Is/I-Am discovery.
  \item \textbf{DNP3 (TCP/20000):} Utility industry standard.
  \item \textbf{S7 (TCP/102):} Siemens proprietary protocol.
\end{itemize}

\section{Methodology}

Our pipeline consists of four stages:

\begin{enumerate}
  \item \textbf{Discovery:} Shodan queries for known ICS ports
  \item \textbf{Enrichment:} Banner analysis and fingerprinting
  \item \textbf{Validation:} Non-intrusive connectivity checks
  \item \textbf{Analysis:} Geographic and sector mapping
\end{enumerate}

\section{Results}

Preliminary analysis reveals exposure patterns across critical infrastructure 
sectors. Specific findings have been redacted pending responsible disclosure.

\section{Ethical Considerations}

We adhere to: (1) Passive-only reconnaissance; (2) Minimal network footprint; 
(3) Structured reporting to US-CERT.

\section{Conclusion}

This paper presents a replicable OSINT methodology for ICS security research.

% For arXiv, include .bbl generated from:
% pdflatex > bibtex > pdflatex > pdflatex
\bibliography{bibliography/references}

\end{document}
